% Figure: the survival-space schematic for the crashworthiness case study.
% The crush stroke of the front structure and the ride-down travel of the
% restrained occupant share the closing distance; the occupant cell must not
% intrude. Drawn as a side schematic with the energy shares annotated.
\begin{figure}[htbp]
\centering
\begin{tikzpicture}[
  cell/.style={draw=engblack, very thick},
  crush/.style={draw=engblue, very thick, fill=engblue!12},
  occ/.style={draw=engamber, very thick},
  arrow/.style={draw=engred, -{Stealth}, very thick},
  dimline/.style={draw=enggray, {Stealth}-{Stealth}, thin},
  label/.style={font=\small},
]
% ground line
\draw[draw=enggray, thin] (-0.3,0) -- (12.3,0);

% the crush zone (front structure) as a corrugated/folding block on the left
\fill[engblue!12] (0,0.2) rectangle (4,2.4);
\draw[crush] (0,0.2) rectangle (4,2.4);
% fold lines suggesting the accordion crush
\foreach \x in {0.7,1.4,2.1,2.8,3.5} {
  \draw[draw=engblue, thin] (\x,0.2) -- (\x-0.25,1.3) -- (\x,2.4);
}
\node[engblue, label, anchor=south] at (2,2.5) {crush zone (stroke $s_c$)};

% the occupant cell (rigid survival space) on the right
\draw[cell] (4,0.2) rectangle (9.5,2.8);
\node[engblack, label, anchor=south] at (6.75,2.85) {occupant cell (must stay rigid)};

% the occupant as a circle (head) + torso, with restraint ride-down travel
\draw[occ] (7.0,1.1) circle (0.45);                 % head
\draw[occ] (7.0,0.65) -- (7.0,0.2);                 % torso to floor
\node[engamber, font=\scriptsize, anchor=west] at (7.5,1.1) {occupant};
% restraint webbing to the cell front
\draw[occ, dashed] (6.6,0.9) -- (4.3,1.7);
% ride-down travel of the occupant inside the cell
\draw[arrow] (7.0,2.0) -- (6.0,2.0);
\node[engamber, font=\scriptsize, anchor=south] at (6.5,2.05) {ride-down $s_r$};

% closing/approach arrow (the vehicle into the barrier)
\draw[arrow] (10.6,1.5) -- (9.8,1.5);
\node[engred, font=\scriptsize, anchor=west] at (10.7,1.5) {$v_0$};

% rigid barrier on the far left
\draw[draw=engblack, very thick] (-0.25,0.1) -- (-0.25,2.6);
\foreach \y in {0.3,0.7,1.1,1.5,1.9,2.3} {
  \draw[draw=enggray, thin] (-0.25,\y) -- (-0.55,\y+0.18);
}
\node[font=\scriptsize, anchor=south, rotate=90] at (-0.62,1.35) {rigid barrier};

% dimension lines for the strokes
\draw[dimline] (0,-0.35) -- (4,-0.35);
\node[enggray, font=\scriptsize, anchor=north] at (2,-0.37) {$s_c$ (structural crush)};
\draw[dimline] (6.0,-0.35) -- (7.0,-0.35);
\node[enggray, font=\scriptsize, anchor=north] at (6.5,-0.55) {$s_r$};

\end{tikzpicture}
\caption[Crash survival space: crush stroke shared with occupant
ride-down]{The closing distance in a frontal impact is shared between two
travels. The front structure crushes through a stroke $s_c$ (left, the
folding blue block), absorbing the bulk of the kinetic energy at a
controlled mean force. The restrained occupant then rides down through a
further travel $s_r$ relative to the cell as the belt and airbag stretch
and pay out, absorbing the occupant's own kinetic energy at a force the
body can survive. The occupant cell between them must stay rigid: any
intrusion into the survival space subtracts directly from $s_r$ and raises
the occupant's deceleration. The design problem is to size $s_c$ and the
mean crush force so the cell never collapses and the occupant's
deceleration over $s_c + s_r$ stays below the survivable limit.}
\label{fig:vol03ch05:crumple-survival-space}
\end{figure}
